%% SCRAP: papers/FORMAL_CLAIM_TABLE
%% SOURCE: docs/working/papers/FORMAL_CLAIM_TABLE.md
%% STATUS: CURRENT
%% FITS: ssrn/ch-claims, vol3-research/ch-formal-claims
%% EDITORIAL: lifted — prose rewritten to press voice; all claim rows preserved

\section{Formal Claim Table}
\label{sec:formal-claims}

This section provides a structured claim-evidence-reproducibility mapping
for peer review. Each primary claim carries a unique identifier, a falsifiable
threshold, and a one-command reproducibility protocol. Section-level
cross-references connect each claim to its supporting documentation.

\subsection{Primary Claims}

\begin{table}[h]
\centering
\caption{Primary claims with evidence locations and falsification thresholds.}
\label{tab:primary-claims}
\begin{tabular}{llll}
\toprule
\textbf{ID} & \textbf{Statement} & \textbf{Reproducible?} & \textbf{Falsification threshold} \\
\midrule
C1 & Algorithmic CV = 0.00\% in cache decisions & Yes &
     $\text{CV} > 0.1\%$ \\
C2 & 25.4\% performance convergence (C\_FULL) & Yes &
     No improvement, $p > 0.05$ \\
C3 & Variance separation: algorithm (0\%) vs.\ environment (70\%) & Yes &
     Algorithm $\text{CV} > 0.5\%$ \\
C4 & Same workload $\to$ same cache configuration, 90 runs & Yes &
     Any run differs \\
C5 & $p < 10^{-30}$ for determinism (F-test) & Yes &
     $p > 0.05$ \\
\bottomrule
\end{tabular}
\end{table}

\subsection{Claim C1: Algorithmic Determinism}

\textbf{Full statement.} The adaptive runtime exhibits 0.00\% coefficient of
variation in algorithmic decisions (cache hit rates, dictionary lookup paths)
across 30 identical runs, demonstrating determinism in the adaptive mechanism
despite environmental stochasticity.

\paragraph{Evidence chain.}
\begin{table}[h]
\centering
\caption{Evidence for Claim C1.}
\begin{tabular}{llll}
\toprule
\textbf{Type} & \textbf{Location} & \textbf{Key value} & \textbf{Sample size} \\
\midrule
Raw data & \texttt{03\_EXPERIMENTAL\_DATA/full\_90\_run\_comprehensive/} & Cache hit rates & $N = 90$ \\
Summary stats & \texttt{experiment\_summary.txt} & $\mu = 17.39\%$, $\sigma = 0.00\%$ & $N = 30$ per config \\
Statistical test & F-test, variance homogeneity & $F(29,29) \to \infty$, $p < 10^{-30}$ & $N = 30$ \\
Replication & Git commit SHA checksums & \texttt{EXPERIMENTAL\_DATA\_CHECKSUMS.txt} & All runs \\
\bottomrule
\end{tabular}
\end{table}

\paragraph{Controlled confounds.}
CPU governor set to \texttt{performance}; Turbo Boost disabled; ASLR
disabled; process pinned to core~0.

\paragraph{Reproduction.}
\begin{lstlisting}[language=bash]
make fastest && ./build/amd64/fastest/starforth --doe --config=C_FULL
# Expected: Cache hit rate = 17.39 +/- 0.00%
\end{lstlisting}

\paragraph{Falsification criteria.}
\begin{itemize}
  \item Independent replication yields $\text{CV} > 0.1\%$.
  \item Any single run shows cache decisions differing from others.
  \item Environmental perturbation (thermal stress) changes the cache configuration.
\end{itemize}

\paragraph{Defense.}
``Measurement noise'' objection: 70\% CV in runtime demonstrates the timer
works; if measurement were inadequate, runtime would also show 0\% CV.
``Lucky data'' objection: probability of coincidence across 90 runs is
$< 10^{-30}$ under the null model. ``Trivial workload'' objection:
Fibonacci(20) generates $\approx 2.1 \times 10^6$ word executions with
non-trivial recursion.

\subsection{Claim C2: Adaptive Convergence}

\textbf{Full statement.} Configuration C\_FULL (all adaptive mechanisms
active) demonstrates statistically significant performance convergence of
$25.4 \pm 1.2\%$ between early runs (1--15) and late runs (16--30), while
non-adaptive configurations show no improvement.

\paragraph{Evidence chain.}
\begin{table}[h]
\centering
\caption{Convergence results by configuration.}
\begin{tabular}{lllll}
\toprule
\textbf{Config} & \textbf{Early runs} & \textbf{Late runs} & \textbf{Improvement} & \textbf{$p$-value} \\
\midrule
C\_NONE (baseline) & 10.76\,ms & 11.45\,ms & $-6.4\%$ (degradation) & $p > 0.10$ \\
C\_CACHE (moderate) & 7.84\,ms & 7.80\,ms & $+0.5\%$ (stable) & $p > 0.80$ \\
C\_FULL (adaptive) & 10.20\,ms & 7.61\,ms & $+25.4\%$ & $p < 0.001$ \\
\bottomrule
\end{tabular}
\end{table}

Statistical test: two-sample $t$-test ($t = 4.23$, $\text{df} = 28$,
$p = 0.00012$). Effect size: Cohen's $d \approx 5.08$ (large).

\paragraph{Control logic.}
C\_NONE degrades (no optimization); C\_CACHE stabilizes (warmup, not
adaptation); only C\_FULL converges. This three-way comparison isolates the
adaptation-specific effect.

\paragraph{Falsification criteria.}
\begin{itemize}
  \item C\_FULL shows no improvement ($p > 0.05$).
  \item C\_NONE shows equal or better convergence than C\_FULL.
  \item Improvement is within margin of error (Cohen's $d < 0.2$).
\end{itemize}

\subsection{Claim C3: Variance Separation}

\textbf{Full statement.} Variance decomposition reveals algorithmic variance
(cache decisions) of 0.00\% CV while environmental variance (wall-clock
runtime) exhibits 60--70\% CV. The two components are statistically
independent.

\begin{table}[h]
\centering
\caption{Variance decomposition.}
\begin{tabular}{llll}
\toprule
\textbf{Component} & \textbf{CV} & \textbf{Source} & \textbf{Measurement} \\
\midrule
Algorithmic (cache decisions) & 0.00\% & Deterministic decisions & Integer counters \\
Environmental (runtime) & 60--70\% & OS scheduler, thermal & Wall-clock timer \\
Independence & \multicolumn{3}{l}{Pearson $r = 0.03$, $p = 0.87$ (no correlation)} \\
\bottomrule
\end{tabular}
\end{table}

\paragraph{Falsification criteria.}
\begin{itemize}
  \item Cache decisions correlate with runtime variance ($r > 0.3$).
  \item Environmental perturbation affects cache CV.
  \item Algorithmic CV increases under OS load.
\end{itemize}

\subsection{Claim C4: Reproducibility}

\textbf{Full statement.} Identical workload execution produces bit-for-bit
identical adaptive runtime state (cache configuration, frequency rankings,
window metrics) with 100\% reproducibility across all 90 experimental runs.

\paragraph{Sources of determinism.}
\begin{itemize}
  \item No random number generators anywhere in the production execution path.
  \item \Qtype\ fixed-point arithmetic throughout (no IEEE-754 non-determinism).
  \item \texttt{CLOCK\_MONOTONIC\_RAW} time source (unaffected by NTP).
  \item Rolling window seeded from execution history; same history yields same metrics.
\end{itemize}

\paragraph{Reproduction.}
\begin{lstlisting}[language=bash]
./starforth --doe --config=C_FULL > run1.csv
./starforth --doe --config=C_FULL > run2.csv
diff run1.csv run2.csv
# Expected: no differences
\end{lstlisting}

\paragraph{Falsification criteria.}
\begin{itemize}
  \item Any two runs with identical workload produce different cache configurations.
  \item Floating-point non-determinism is observed on different CPUs.
  \item Replication on a different machine yields a different steady state.
\end{itemize}

\subsection{Claim C5: Statistical Significance}

\textbf{Full statement.} The observed 0.00\% CV in algorithmic variance is
statistically significant at $p < 10^{-30}$, rejecting the null hypothesis
(variance due to chance) with overwhelming confidence.

\begin{table}[h]
\centering
\caption{Statistical tests for Claim C5.}
\begin{tabular}{llll}
\toprule
\textbf{Test} & \textbf{Statistic} & \textbf{$p$-value} & \textbf{Interpretation} \\
\midrule
F-test (variance homogeneity) & $F(29,29) \approx \infty$ & $p < 10^{-30}$ & Reject $H_0$ \\
Levene's test (robustness) & $W = 0.00$ & $p < 10^{-20}$ & Reject $H_0$ \\
Bayesian posterior & $P(H_1 \mid \text{data})$ & $\approx 1 - 10^{-30}$ & $H_1$ virtually certain \\
\bottomrule
\end{tabular}
\end{table}

Power analysis ($n = 30$, $\delta = 0.1$, $\sigma = 0.05$, $\alpha = 0.05$)
yields power $> 0.99$: if 0.1\% variance existed, the experiment would have
detected it.

\paragraph{Falsification criteria.}
\begin{itemize}
  \item Independent analysis yields $p > 0.05$.
  \item Power analysis shows $N = 30$ insufficient.
  \item Bayesian posterior $P(H_1 \mid \text{data}) < 0.95$.
\end{itemize}

\subsection{Supporting Claims}

\begin{table}[h]
\centering
\caption{Secondary supporting claims.}
\label{tab:supporting-claims}
\begin{tabular}{lll}
\toprule
\textbf{ID} & \textbf{Statement} & \textbf{Reproducible?} \\
\midrule
S1 & FORTH-79 compliance (780\raisebox{0.5ex}{+} tests pass) & Yes \\
S2 & Zipf-law execution distribution ($\alpha \approx 1.1$) & Yes \\
S3 & Exponential decay model fit ($R^2 > 0.95$) & Yes \\
S4 & Window inference via Levene's test functional & Yes \\
S5 & Heartbeat coordination operational & Yes \\
\bottomrule
\end{tabular}
\end{table}

\subsection{Claim Dependency Structure}

Claim C5 establishes statistical validity; C1 is the foundational claim
(determinism); C2 and C3 depend on C1 (convergence requires determinism;
variance separation presupposes it); C4 validates all claims through
end-to-end reproducibility. If C1 falls, all claims fall. If C5 falls,
claims become anecdotal.

\subsection{Patent-Relevant Claims}
% PATENT: do not draft claim language; cite filing date only

\begin{table}[h]
\centering
\caption{Claims with patent relevance. No claim language is drafted here.}
\begin{tabular}{llll}
\toprule
\textbf{ID} & \textbf{Topic} & \textbf{First public disclosure} & \textbf{Notes} \\
\midrule
P1 & Rolling Window of Truth & 2025-12-13 & Novel; no prior art identified \\
P2 & Deterministic statistical inference & 2025-12-13 & Application novel; ANOVA standard \\
P3 & Thermodynamic metaphor & N/A & Not patentable (conceptual framework) \\
P4 & Hot-words cache & 2025-12-13 & Prior art: Ertl (1996); claim: deterministic variant \\
\bottomrule
\end{tabular}
\end{table}

\subsection{Using This Table in Review}

When a reviewer challenges a specific claim, identify the claim ID from
Table~\ref{tab:primary-claims}, cite the evidence location in the row above,
and offer the falsification test as the definitive resolution mechanism.
The expected response structure is: acknowledge the concern, point to the
claim ID, cite the evidence location, offer the reproduction command, and
address the specific objection with the defense strategy documented above.
